\documentclass[12pt]{article}
\usepackage[a4paper,margin=2.5cm]{geometry}
\usepackage[T1]{fontenc}
\usepackage[utf8]{inputenc}
\usepackage{lmodern}
\usepackage{microtype}
\usepackage{xcolor}
\usepackage{listings}
\usepackage[hidelinks]{hyperref}
\definecolor{codegray}{RGB}{245,245,245}
\lstset{basicstyle=\ttfamily\small,backgroundcolor=\color{codegray},frame=single,
  breaklines=true,showstringspaces=false,columns=fullflexible,keepspaces=true}
\setlength{\parindent}{0pt}
\setlength{\parskip}{0.55em}

\begin{document}
\pagenumbering{gobble}
\begin{titlepage}
\centering
{\large Ministerul Educa\c{t}iei \c{s}i Cercet\u{a}rii al Republicii Moldova\par}
\vspace{0.45cm}
{\large Universitatea Tehnic\u{a} a Moldovei\par}
\vspace{0.45cm}
{\large Facultatea Calculatoare, Informatic\u{a} \c{s}i Microelectronic\u{a}\par}
\vspace{0.45cm}
{\large Departamentul Ingineria Software \c{s}i Automatic\u{a}\par}
\vspace{3cm}
{\Large\bfseries Report on Laboratory Work No. 7\par}
\vspace{0.55cm}
{\LARGE\bfseries Processes, Redirection, and Pipes\par}
\vspace{0.35cm}
{\large Operating Systems\par}
\vfill
\begin{flushleft}
\textbf{Done by:} st. gr. FAF-241, Ca\c{t} Chiril\\[0.8cm]
\textbf{Checked by:} asist. univ. Rei\c{t}man Patricia
\end{flushleft}
\vfill
{Chi\c{s}in\u{a}u -- 2026\par}
\end{titlepage}

\pagenumbering{roman}
\tableofcontents
\clearpage
\pagenumbering{arabic}

\section{Purpose and environment}
I built and ran four POSIX C programs to observe process creation, program
replacement, child reaping, standard-output redirection, and a two-stage pipe.
The experiments ran on macOS on 26 October 2026. The \texttt{gcc} command
reported Apple clang 21.0.0. I compiled each source with
\texttt{-std=c11 -Wall -Wextra -Werror} and used a temporary directory for
binaries and generated files. The submitted \texttt{code/} directory contains
only the four required C sources. A second source copy is stored in
\texttt{xv6-riscv/lab7-linux/}; the programs use POSIX APIs and are outside
xv6's \texttt{UPROGS}. The required \texttt{report.md} contains the full
Observe answers and selected command outputs.

\section{Fork, exec, and wait}
The final \texttt{run.c} forks, executes \texttt{ls -l} in the child, and
waits in the parent. For the two requested variations, I compiled temporary
copies that replaced the \texttt{execlp} line with \texttt{date} and
\texttt{nosuchprog}; the submitted source remains the original \texttt{ls}
version.

\begin{lstlisting}
$ ./run
parent: my PID is 55082
child:  my PID is 55084, my parent is 55082
total 536
... (directory listing)
parent: child 55084 ended with status 0
$ ./run_date
parent: my PID is 52683
child:  my PID is 52695, my parent is 52683
Mon Oct 26 11:05:35 EET 2026
parent: child 52695 ended with status 0
$ ./run_missing
parent: my PID is 52681
child:  my PID is 52694, my parent is 52681
exec: No such file or directory
parent: child 52694 ended with status 127
\end{lstlisting}

\texttt{perror("exec")} runs only for \texttt{nosuchprog} because a successful
\texttt{execlp} replaces the child's process image and never returns. The
child's parent value is the original process's PID: 55082 in the \texttt{ls}
run. The parent observes status 0 for \texttt{ls} and 127 for the missing
program, matching the child's explicit \texttt{exit(127)} after the error.

\clearpage
\section{A zombie before waitpid}
\texttt{zombie.c} is the original \texttt{ls} program with
\texttt{sleep(20)} inserted before \texttt{waitpid}. After two seconds, the
parent was sleeping and its child had exited but had not yet been reaped:

\begin{lstlisting}
$ ./zombie > /dev/null &
$ sleep 2
$ ps -o pid,ppid,stat,command -p 54291
  PID  PPID STAT COMMAND
54291 54289 SN   ./zombie
$ ps -axo pid,ppid,stat,command | awk '$2 == 54291 {print}'
54293 54291 ZN   <defunct>
$ wait
$ ps -axo pid,ppid,stat,command | awk '$2 == 54291 {print}'
(no rows)
\end{lstlisting}

The \texttt{Z} in \texttt{STAT} means zombie. The additional \texttt{N} is a
macOS scheduling flag. After 20 seconds the defunct child disappeared because
the parent's \texttt{waitpid(pid, \&status, 0)} collected its exit status.

\section{Redirecting standard output}
The child in \texttt{redirect.c} closes descriptor 1, opens \texttt{out.txt},
then executes \texttt{ls -l}. The new file gets the lowest available number:

\begin{lstlisting}
$ ./redirect
open gave descriptor 1
done: the listing is in out.txt
$ head -3 out.txt
total 448
-rw-r--r--@ 1 chirilcat  wheel      0 Oct 26 11:05 out.txt
-rwxr-xr-x@ 1 chirilcat  wheel  33848 Oct 26 11:05 pipeline
\end{lstlisting}

\texttt{open} returned descriptor 1 because it was just closed. The unchanged
\texttt{ls} writes to descriptor 1, which now refers to the file, so its
listing is absent from the screen. The diagnostic still goes to descriptor 2
(standard error). Redirecting errors would require closing descriptor 2
before opening the error file, or assigning the file to 2 with \texttt{dup2}.

\clearpage
\section{Connecting commands with a pipe}
\texttt{pipeline.c} uses a pipe and two forks. The \texttt{ls} child maps the
pipe's writing end to descriptor 1; the \texttt{wc} child maps its reading
end to descriptor 0. Each process closes pipe descriptors it no longer uses.

\begin{lstlisting}
$ ./pipeline
      77
$ ls /etc | wc -l
      77
\end{lstlisting}

The C program and the shell pipeline produced the same count. I then built
a temporary version with the parent's \texttt{close(p[1]);} commented out.
After two seconds, \texttt{ls} had exited while \texttt{wc} and the parent
were still waiting:

\begin{lstlisting}
  PID  PPID STAT COMMAND
54803 54801 SN   ./pipeline_broken
54806 54803 SN   wc -l
\end{lstlisting}

\texttt{wc} cannot observe end-of-file while the parent still holds a copy
of the pipe's write end. I stopped the temporary processes after recording
their state. The submitted \texttt{pipeline.c} closes both pipe ends in the
parent.

\section{Conclusion}
\texttt{fork} creates a child, \texttt{exec} replaces its program, and
\texttt{waitpid} collects its exit status. Redirection and pipelines assign
files or pipe endpoints to standard descriptors before \texttt{exec}. The
broken-pipe experiment showed why every unused write end must be closed so
the reader can receive end-of-file.

\end{document}
